\section{Introduction}

Les expériences ont été réalisées sur un pendule de Pohl. Il s'agit d'un oscillateur de torsion composé d’un disque monté sur un axe et rappelé par un ressort spiral. Nous avons experimente differents régimes d’oscillation sur ce dernier.
\\ \\
Dans la première expérience, nous avons observé les oscillations libres amorties par un frein électromagnétique. Dans la seconde expérience, le pendule a été soumis à une excitation forcée par un moteur.
\\ \\
L’objectif de ce travail est de caractériser le comportement dynamique du système. A partir des mesures des oscillations, il a été possible de déterminer les paramètres physiques du pendule tels que la pulsation propre et le coefficient d’amortissement. Cela nous a permit de vérifier expérimentalement certaines lois fondamentales, comme la loi de Hooke, le moment cinétique ou encore l’effet du frein de Foucault dans un regime d'amortissement faible. L'analyse du systeme nous a egalement permi de determiner constante de torsion du ressort spiral utilisé.

\section{Description de l'installation}

En premier lieu, voici l'installation du pendule de pohl dont on peut clairement voir le balancier et le ressort spiral que l'on va manipuler durant les deux experiences

\\
\begin{figure}[ht]
\centering
\includegraphics[scale=0.8]{schema_pohl.png}
	\caption{ Schema de Pohl, voir ref : \cite{schemapohl} }
\label{fig:schemapohl}
\end{figure}

Comme il est possible de le voir en fig:\ref{fig:schemapohl}, le dispositif est constitué d’une roue montée sur un axe et reliée à un ressort spiral, dont l’autre extrémité est fixée ou mise en mouvement par un systeme de transmission liant le balancier au moteur. La roue du balancier passe dans l’entrefer d’un frein de Foucault. Le frein permet de moduler l’amortissement en ajustant le courant mesure par un amperemetre (voir fig :\ref{samplesetup}). Le mouvement angulaire du pendule est enregistré par un encodeur connecté à l’interface Phidgets MATLAB. Cela elimine grandement les erreurs accidentelles de mesures, car les mesures sont principalement electroniques. Des indexes sur le pendule et l’excitateur permettent de mesurer la phase relative.

\newline 
L'ensemble* des experiences sont realisables grace au schema ci-dessous : 
\newline
\begin{figure}[ht]
\centering
\includegraphics[scale=0.5]{schema_force.png}
\caption{Schema electrique pour l'analyse du pendule de Pohl}
\label{fig:samplesetup}
\end{figure}

\\
Le schema est strictement le meme pour les experiences en regime libre et en regime force. Seul le moteur est branche pour la seconde experience. Lorsqu'il sera alimente, cela modifiera l'quation caracteristique de notre systeme dynamique (voir Eq. \ref{eq:cara}) afin de rajouter un terme a droite de l'equaion (voir Eq. \ref{eq:equa_forcee_normalisee}.
\newline \\
*Il manque seulement aux schemas le dynamometre qui a permit de mesurer la constante de torsion C du ressort-spiral (voir Eq. \ref{eq:C}). Avant de faire osciller le systeme, il suffit de placer suspendre un dynamometre au dessus du pendule et de le fixer au moyeu (cylindre du centre de rotation) de ce dernier. Afin de realiser les mesures pour determiner C, il faut mesurer le rapport du moment de force avec le dynamometre et l'angle responsable de ce moment de force. La suite des mesures est moins manuelle.

\newline \\ \\
Pour les oscillations libres, le disque est dévié d’un angle initial et lâche. Le ressort spiral exerce un moment de rappel qui oppose la rotation et le pendule effectue des oscillations amorties. L’amplitude et la pseudo-période sont mesurées pour différentes valeurs de courant dans le frein afin de déterminer le coefficient d’amortissement et le moment d’inertie. 
\newline \\
Pour les oscillations forcées, le moteur impose un mouvement sinusoïdal de fréquence réglable. L’amplitude et la phase de l’oscillateur sont enregistrées en fonction de la fréquence, permettant de tracer les courbes caractéristiques et d’identifier la fréquence de résonance ainsi que le facteur de qualité.
\newline\\
Des mesures complémentaires sont réalisées avec un ampermètre pour quantifier le courant induit dans le frein et relier ainsi l’amplitude des oscillations à l’intensité du freinage. L’ensemble du montage est simple à manipuler tout en offrant un contrôle précis des paramètres dynamiques du pendule.


